\section{Realization in a central division algebra}
\label{sec:division-realization}

The preceding construction solved equations in split matrices.  We now
place its constant blocks and prescribed arrows in one central division
algebra.  Each subsequent correction is the solution of a finite linear
system over its center.  This permits descent without any assertion that
solutions of nonlinear equations descend.

Throughout this section $k$ has characteristic zero and contains all roots
of unity, and $A$ is a central division algebra of degree $d$ over $k$.
The integer $b$ is the block size chosen earlier in the spectral case.
In the case where $q_0(w)=0$ already holds as a free-series identity, any
positive integer $b$ may be used.  Put $m=11$.

\subsection{The two division-algebra factors}

Choose independent central indeterminates $\alpha,\beta$ and put
$k_0=k(\alpha,\beta)$.  Let $B_b$ be the symbol algebra over $k_0$
generated by $u,v$ with
\[
 u^b=\alpha,\qquad v^b=\beta,\qquad vu=\zeta_buv,
\]
where $\zeta_b$ is a primitive $b$-th root of unity.  For $b=1$ this
means $B_1=k_0$.  By the construction in Section~\ref{subsec:generic-division}, $B_b$
is central simple of degree $b$.  Lemma~\ref{lem:symbol-division}
shows that $A\otimes_k B_b$ is division and remains division under
pure rational and central Laurent scalar extensions.  We will use
these facts twice in adjoining the cyclic block factor.

Fix a $k_0$-basis $(e_a)_{a=1}^{b^2}$ of $B_b$.  Over $k_0$ introduce
the following jointly independent variables:
\[
 x_{c,a,j},\qquad y_{r,c,a,j},
 \qquad
 \begin{matrix}
 c\in\{1,2\},\quad 1\le a\le b^2,\\
 j\in\mathbb Z/m\mathbb Z,\quad r\in\{1,2,3,4\}.
 \end{matrix}
\]
Let $L_0$ be their rational function field and let $\sigma$ increase
the index $j$ by one, fixing $k_0$.  This automorphism has order exactly
$m$.  Put $H_0=L_0^{\langle\sigma\rangle}$.  Artin's fixed-field theorem
for a finite group of automorphisms gives a cyclic Galois extension
$L_0/H_0$ of degree $m$; see \cite[Corollary~3.11, p.~38]{MilneFT2022}.

Choose a further independent central indeterminate $a_*$ and form the
cyclic algebra $C$ over $F_0=H_0(a_*)$ with coefficient field
$L_0(a_*)$ and generator $\tau$, subject to
\begin{equation}
 \tau^m=a_*,\qquad \tau c\tau^{-1}=\sigma(c)
 \quad(c\in L_0).
 \label{eq:cyclic-realization-relations}
\end{equation}
It has vector-space decomposition
$C=\bigoplus_{r=0}^{m-1}L_0(a_*)\tau^r$.
After extending scalars to a field containing $L_0(a_*)$, represent
$c\in L_0$ by the diagonal with entries $\sigma^{-i}(c)$ and
$\tau$ by the cyclic forward shift, whose wrap entry is $a_*$.
These matrices satisfy \eqref{eq:cyclic-realization-relations} and
span all matrix units.  This gives the usual splitting of $C$ and
shows that it is central simple of degree $m$; see
\cite[Construction~2.5.1 and Proposition~2.5.2]{GilleSzamuely2006}.
Our conjugation convention uses the inverse of their Galois generator.

Define
\[
 \widetilde B_0=(B_b\otimes_{k_0}F_0)\otimes_{F_0}C,
 \qquad D_0=A\otimes_k\widetilde B_0.
\]
Then $D_0$ is central simple of degree $dbm$ over $F_0$ by the tensor-product and scalar-extension properties of
central simple algebras
\cite[Chapter~IV, Corollary~2.8 and Proposition~2.15]{MilneCFT}.
It is also division. Indeed
\[
 D_L=(A\otimes_k B_b)\otimes_{k_0}L_0
\]
is division by the rational-extension argument.  Extend $\sigma$ to
$D_L$ by fixing the old algebra factors.  The skew Laurent division
ring $D_L((t;\sigma))$ receives a unital homomorphism from $D_0$
by $\tau\mapsto t$ and $a_*\mapsto t^m$.  The central element
$t^m$ is transcendental over $H_0$.  Simplicity again makes this map
injective, and finite dimension makes $D_0$ division.

Finally put
\[
 F=F_0((\varepsilon)),\qquad
 \widetilde B=\widetilde B_0\otimes_{F_0}F,\qquad
 D=D_0\otimes_{F_0}F.
\]
The central Laurent argument shows that $D$ is a central division
algebra over $F$ of degree $dbm$.  Input coefficients will lie in
$\widetilde B$, so that they commute with the old factor $A$.

\subsection{One realization of all generic data}

For $r\in\mathbb Z/m\mathbb Z$, let
\[
 \mathcal V_r=
 \bigl(B_b\otimes_{k_0}L_0\bigr)\tau^r
 \otimes_{H_0}F
 \ \subseteq\ \widetilde B.
\]
This notation uses the evident coefficient embedding and scalar
extension; equivalently it is the $F$-span of that graded summand in
$\widetilde B_0$.  Upon splitting, $\mathcal V_r$ consists exactly
of all $b\times b$ entries at offset $r$ in the $m$-block matrix.
In fact $L_0\otimes_{H_0}E\simeq E^m$ over a splitting field $E$,
and $B_b$ becomes $M_b(E)$.  Multiplication by $\tau^r$ moves these
independent diagonal entries to offset $r$, with nonzero wrap factors.
In particular $\mathcal V_0$ gives all diagonal blocks, and
$\mathcal V_{-r}$ gives all reverse entries for offset $r$.

Define the two constant inputs and four arrow pairs by
\[
 P_c^0=\varepsilon\sum_a x_{c,a,0}e_a,
 \qquad
 U_{r,c}=\left(\sum_a y_{r,c,a,0}e_a\right)\tau^r
 \quad(c=1,2).
\]
The prescribed initial substitution is the pair
\begin{equation}
 Z_c^{\mathrm{init}}(s_1,s_2)
 =P_c^0+(U_{1,c}+U_{2,c})s_1
             +(U_{3,c}+U_{4,c})s_2.
 \label{eq:division-initial-input}
\end{equation}
After splitting, the diagonal blocks of $(P_1^0,P_2^0)$ are exactly
the required independent small generic pairs.  The two input-component
entries at every forward edge are independent generic entries, jointly
with all the diagonal data.

Here independence is over the old split coefficient field, which is the
field used in the generic rank tests.  To check it, first split $A$ and
$B_b$ using constants independent of the variables $x,y$.  The $m$
conjugates of each coefficient use the disjoint sets of variables with
different index $j$.  The split symbol basis is a matrix basis.
Permuting those variables, changing that basis, and multiplying entries
by powers of $a_*$ preserve their algebraic independence.  This verifies
all generic determinant conditions simultaneously.  No independence
over the larger invariant field $H_0$ is asserted or needed.

We may perform the spectral calculations after extending $F$ to a
split Laurent field and then to the finite extensions needed for
eigenvalues.  Its constant coefficient field contains the old algebraic
closure and the rational variables just chosen.  The matrix constant
terms remain $\varepsilon$-small.  Arrow coefficients need not be
small, since they occur only with positive free-word degree.

\subsection{Linear descent of the correction}

Assume first that $q_0(w)\ne0$.  Let $W_0=w(P_1^0,P_2^0)\in D$
and choose its monic minimal polynomial
\begin{equation}
 q(t)\in F[t].
 \label{eq:fixed-central-polynomial}
\end{equation}
We keep this polynomial fixed throughout the correction.  It is
irreducible, because $F[W_0]$ is a finite-dimensional commutative
domain, and separable because the characteristic is zero.  In a
splitting field its roots are exactly the distinct eigenvalues of the
constant output.  For completeness, left multiplication by $W_0$ on
$D$ has characteristic polynomial a power of $q$, since $D$ is a
vector space over $F[W_0]$.  After splitting, that left multiplication
is a direct sum of copies of the action of the split matrix $W_0$.
Every root of $q$ therefore occurs in that matrix, and no other
eigenvalue occurs.  Separability makes its action semisimple.

We solve $q(w(Z))=0$ in free series over $D$.  Initially take
\eqref{eq:division-initial-input}.  Its split Schur equations vanish
through degree two, so Lemma~\ref{lem:schur-polynomial} gives the same
assertion for the polynomial relation.  Suppose inductively that a
partial substitution solves this relation through degree $\ell-1$,
where $\ell\ge3$.  Allow a new coefficient in
$\mathcal V_{-r}^2$ at every word of length $\ell-1$, for
$r=1,2,3,4$, and a new coefficient in $\mathcal V_0^2$ at every word
of length $\ell$.  These are finite-dimensional $F$-spaces.

Impose every coefficient equation of
\begin{equation}
 q(w(Z))=0\pmod{(s_1,s_2)^{\ell+1}}.
 \label{eq:finite-central-correction}
\end{equation}
This is a finite affine linear system over $F$ in the allowed new
coordinates.  Indeed two new insertions have degree at least
$2(\ell-1)>\ell$.  At any fixed free-word degree, only boundedly
many positive-degree insertions occur in the evaluation of the old
series; the remaining insertions are from the $\varepsilon$-small
constant input.  The finitely many previously fixed coefficients have
a finite lower valuation bound.  Thus the sums defining each
coefficient of \eqref{eq:finite-central-correction}, and its linear
part in the new coordinates, converge in $D$ and are defined over $F$.
Choosing an $F$-basis of $D$ turns them into finitely many ordinary
scalar linear equations.

After scalar extension, each permitted input space becomes exactly
the collection of reverse edges or loops used in
Proposition~\ref{prop:eleven-block-correction}.  The current partial
substitution is a valid partial split solution by
Lemma~\ref{lem:schur-polynomial}.  The proposition supplies a split
correction, and that lemma converts its Schur vanishings back to all
the equations in \eqref{eq:finite-central-correction}.  In particular
the new degree-$\ell-1$ reverse coefficients preserve the already
solved equations: their diagonal Schur variation at that degree is
zero.  We have therefore solved the whole finite linear system after
a field extension.  A matrix and its augmented matrix have the same
ranks before and after extension, so the system already has a solution
over $F$.

Choose such a solution at every step.  The coefficient stabilization
in Proposition~\ref{prop:eleven-block-correction} produces
\[
 Z\in\widetilde B\langle\!\langle s_1,s_2\rangle\!\rangle^2,
 \qquad q(w(Z))=0,
\]
with the prescribed constant and linear terms.  Only finite linear
systems have been descended.  Neither the roots of $q$ nor the split
correction coefficients need belong to $F$.

Since the input coefficients commute with $A$, evaluation at $Z$
defines a substitution homomorphism
\[
 A\langle\!\langle z_1,z_2\rangle\!\rangle
 \longrightarrow D\langle\!\langle s_1,s_2\rangle\!\rangle.
\]
It is well-defined by the same coefficientwise valuation bound: at
fixed output degree, only boundedly many nonconstant input factors
occur.  The bound also justifies multiplication and regrouping of the
convergent coefficient sums.  Hence every old polynomial identity on
the tuple and the chosen series is preserved by this substitution.

If $q_0(w)=0$ already holds, use the very same division algebra and
the initial substitution \eqref{eq:division-initial-input}, with any
$b\ge1$.  No spectral correction is required; the old relation
$q_0(w)$ remains zero under the substitution homomorphism.  Both cases
produce a central division algebra of degree $11bd>d$, with identical
generic forms for the constant and the two linear input coefficients.
The next step restores generation and two transverse directions using
these coefficients.  That step is needed in the unconstrained case
as well as in the spectral case.
